\documentclass[12pt]{article}
\usepackage[utf8]{inputenc}
\usepackage[T1]{fontenc}
\usepackage[paperwidth=21.587cm,paperheight=27.937cm,
            left=3cm,right=2.5cm,top=2.5cm,bottom=2.5cm]{geometry}
\usepackage{amsmath,amssymb,amsfonts}
\usepackage{xcolor}
\usepackage{hyperref}
\usepackage{float}
\usepackage{booktabs}
\usepackage{array}
\usepackage{multirow}
\usepackage{caption}
\usepackage{subcaption}
\usepackage{graphicx}
\usepackage{setspace}
\onehalfspacing

\definecolor{myblue}{RGB}{0,114,178}
\definecolor{myred}{RGB}{213,94,0}

\hypersetup{colorlinks=true,linkcolor=myblue,citecolor=myblue,urlcolor=myblue}

\title{\textbf{Likelihood Estimation for the Nonlinear Panel Data Framework\\
of Arellano, Blundell and Bonhomme (2017):\\
Discontinuities, Efficiency, and a Smooth Extension}\\[0.6em]
\large Preliminary Draft}
\author{Hidehiko Ichimura and Risheng Xu\\
University of Arizona}
\date{\today}

\begin{document}
\maketitle

\begin{abstract}
\noindent
Arellano, Blundell, and Bonhomme (2017, ABB) propose a nonlinear panel data
framework for earnings and consumption dynamics, estimated via a stochastic EM
algorithm with a quantile-regression (QR) M-step.
This paper investigates three aspects of the ABB framework.
First, the piecewise-uniform density with exponential tails creates
discontinuities at interior quantile nodes, causing difficulties for
likelihood-based estimation.
Second, the QR estimator, while computationally tractable, does not maximise
the likelihood of the observed data and is therefore inefficient under correct
specification.
Third, QR targets quantiles of unobservable components ($\eta$ and
$\varepsilon$), and it is unclear how well these translate to the distribution
of the observables $y$.
In contrast, MLE directly optimises the likelihood of observed $y$.

We make three contributions.
(1) We conduct Monte Carlo comparisons of simulated MLE (SML), grid-based MLE,
and QR under both the original ABB density and a smooth extension.
(2) We propose a smooth version of the ABB model using a $C^2$ cubic spline
log-density with Gaussian tails, which is a direct extension of the
Arellano--Bonhomme (2012) and ABB quantile framework while connecting to the
log-spline density estimation literature.
(3) We analyse sources of difficulty for MLE, including the role of
weak identification of tail parameters through deconvolution, the effect of
data generation accuracy, and the information structure via the Louis (1982)
decomposition.
\end{abstract}

\newpage
\tableofcontents
\newpage

%==========================================================================
\section{Introduction}
\label{sec:intro}
%==========================================================================

Household earnings and consumption dynamics over the life cycle are among the
most studied topics in empirical economics.
Arellano, Blundell, and Bonhomme (2017, hereafter ABB) make a landmark
contribution by developing a fully nonlinear panel data framework in which the
persistent earnings component $\eta_{it}$ follows a first-order Markov process
whose conditional quantile function is a flexible polynomial series in
$(\eta_{i,t-1}, \text{age}_{it})$.
Their framework allows earnings persistence to vary with the sign and magnitude
of the current shock, producing the empirically prominent ``nonlinear
persistence'' pattern documented in PSID and Norwegian administrative data.

ABB estimate the model using a stochastic EM algorithm whose M-step consists of
a sequence of quantile regressions (QR) rather than likelihood maximisation.
This paper examines three limitations of this approach and proposes remedies.

\subsection*{Limitation 1: Discontinuities in the Density}

The ABB parameterisation specifies the conditional density of
$\varepsilon_{it} = y_{it} - \eta_{it}$ as piecewise uniform on the intervals
between quantile knots $A(\ell)$, with exponential tails outside
$[A(1), A(L)]$.
The resulting density is \emph{discontinuous} at the interior knots
$A(2), \ldots, A(L-1)$: it jumps between segments of different heights.
Furthermore, the density of $\eta_{it} \mid \eta_{i,t-1}$ has the same
piecewise-uniform structure, so the conditional density of
$y_{it} \mid \eta_{i,t-1}$ (a convolution) inherits kinks from both sources.

These discontinuities make the log-likelihood non-differentiable, preventing
the use of gradient-based optimisation.
When derivative-free methods (e.g., Nelder--Mead) are used, the flat interior
segments allow parameter vectors far from the truth to attain comparable or
higher likelihood values, so the MLE does not reliably locate the true
parameter.

\subsection*{Limitation 2: Efficiency Loss of the QR Estimator}

The QR M-step minimises the check function
$\rho_\tau(u) = u(\tau - \mathbf{1}\{u \le 0\})$
rather than the log-likelihood.
Under correct specification, the QR estimator is consistent but
\emph{not efficient} in the Cram\'er--Rao sense.
The efficiency loss depends on the density shape between the quantile knots
--- information that QR discards.

\subsection*{Limitation 3: QR Targets Unobservables, Not the Distribution of $y$}

A more fundamental concern is that the QR M-step estimates quantile parameters
of the \emph{unobservable} components $\eta_{it}$ and $\varepsilon_{it}$,
conditional on draws from the posterior $p(\eta \mid y)$.
It does not directly optimise any criterion involving the distribution of the
\emph{observed} data $y_{it}$.
In contrast, MLE maximises the likelihood of $y$, which is a convolution
\begin{equation}\label{eq:conv}
f(y_{it} \mid y_{i,t-1}; \theta) = \int f(\eta_{it} \mid \eta_{i,t-1}; \theta_Q)\,
f_\varepsilon(y_{it} - \eta_{it}; \theta_\varepsilon)\, d\eta_{it}.
\end{equation}
In the parametric model we consider, the quantile parameters
$\theta_Q$ and $\theta_\varepsilon$ jointly determine the density of $y$.
QR estimates these parameters via moment conditions on the unobservables,
while MLE uses the full information in the joint density of $y$.
This distinction becomes particularly relevant under model misspecification:
QR fits quantiles of $\eta$ and $\varepsilon$ well but may not approximate
the distribution of $y$ well, whereas MLE directly targets the density of $y$.

We develop this argument formally in the context of our smooth parametric
model in Section~\ref{sec:identification}.

\subsection*{Our Contributions}

This paper makes three contributions.

\begin{enumerate}
\item \textbf{Monte Carlo comparisons.}
  We compare simulated MLE, grid-based MLE, and QR under both the original
  ABB density and a smooth extension, documenting bias, variance, and RMSE
  across sample sizes and number of time periods.

\item \textbf{Smooth ABB model.}
  We propose a $C^2$ cubic spline log-density with Gaussian (quadratic)
  tails as a direct smooth extension of the AB/ABB quantile framework.
  The model preserves the quantile-knot interpretation of ABB while
  ensuring a smooth, everywhere-positive density with analytical gradients
  and exact normalisation via Taylor series integration.

\item \textbf{Sources of difficulty for MLE.}
  We identify and analyse several sources of difficulty:
  (a) the data generation procedure (coarse-grid CDF vs.\ rejection sampling),
  (b) the gradient computation method (numerical vs.\ analytical),
  (c) the parameterisation of tail curvature (free vs.\ profiled from IQR),
  (d) the information structure via the Louis (1982) decomposition, and
  (e) the connection to deconvolution ill-posedness from the
  nonparametric identification theory of Hu and Schennach (2008) and
  Wilhelm (2015).
\end{enumerate}

The remainder of the paper is organised as follows.
Section~\ref{sec:abb} reviews the ABB model.
Section~\ref{sec:mc_abb} presents Monte Carlo evidence comparing MLE and QR
under the ABB density.
Section~\ref{sec:smooth} introduces the smooth ABB model.
Section~\ref{sec:mc_smooth} reports Monte Carlo evidence for the smooth model.
Section~\ref{sec:identification} analyses identification and information.
Section~\ref{sec:conclusion} concludes.

%==========================================================================
\section{The ABB Model}
\label{sec:abb}
%==========================================================================

\subsection{Model Setup}

We consider a panel of $N$ individuals observed over $T$ periods.
The observed outcome is
\begin{equation}\label{eq:model}
y_{it} = \eta_{it} + \varepsilon_{it}, \quad i=1,\ldots,N, \quad t=1,\ldots,T,
\end{equation}
where $\eta_{it}$ is a persistent component following a first-order Markov
process and $\varepsilon_{it}$ is an i.i.d.\ transitory shock with zero mean,
independent of $\eta_{is}$ for all $s$.

The conditional quantile function of $\eta_{it}$ given $\eta_{i,t-1}$ is
specified as
\begin{equation}\label{eq:quantile}
Q_\tau(\eta_{i,t-1}) = \sum_{k=0}^K a_{k}(\tau)\, H_k(\eta_{i,t-1}/\sigma_y),
\end{equation}
where $H_k$ are Hermite polynomials and $\tau \in \{\tau_1, \ldots, \tau_L\}$
is a grid of quantile levels (e.g., $\tau = 0.25, 0.50, 0.75$ for quartiles).

\subsection{Piecewise-Uniform Density with Exponential Tails}

The density implied by the piecewise-linear quantile function is
piecewise uniform on the intervals $[Q_{\tau_\ell}, Q_{\tau_{\ell+1}})$,
with exponential tails below $Q_{\tau_1}$ and above $Q_{\tau_L}$.
The density of $\varepsilon_{it}$ has the same structure.

This density is \emph{discontinuous} at the interior knots.
Figure~\ref{fig:disc_density} illustrates the density of $\eta_{i1}$
under the ABB parameterisation, showing the jumps at quantile boundaries.

\textcolor{myred}{[TODO: Insert figure showing piecewise-uniform density
with discontinuities, and the implied density of $y$ with kinks.]}

\subsection{Parameter Counting}

With $K=2$ (quadratic Hermite basis) and $L=3$ (quartiles), the model has:
\begin{itemize}
\item Transition: $a_Q \in \mathbb{R}^{(K+1) \times L} = \mathbb{R}^{3 \times 3}$
  (9 parameters) plus 2 tail rates $= 11$.
\item Initial $\eta_1$: 3 quantile knots plus 2 tail rates $= 5$.
\item Transitory $\varepsilon$: 2 quantile knots (median $=0$) plus 2 tail rates $= 4$.
\item Total: $11 + 5 + 4 = 20$ parameters.
\end{itemize}

%==========================================================================
\section{Performance of SML, MLE, and QR in the ABB Model}
\label{sec:mc_abb}
%==========================================================================

\textcolor{myred}{[TODO: This section will contain Monte Carlo comparison results
for the original ABB (piecewise-uniform) model.
Simulation design and results to be completed.
Will include SML, grid-based MLE, and QR comparisons across $N$ and $T$.]}

%==========================================================================
\section{The Smooth ABB Model}
\label{sec:smooth}
%==========================================================================

\subsection{Motivation}

The discontinuities in the ABB density motivate a smooth alternative.
We seek a model that:
\begin{enumerate}
\item preserves the quantile-knot structure of ABB (each inter-quantile
  segment contains exactly $\tau_{\ell+1} - \tau_\ell$ probability mass);
\item produces a $C^2$ (twice continuously differentiable) density,
  enabling gradient-based MLE;
\item has integrable tails with well-defined curvature; and
\item connects naturally to the quantile regression interpretation of AB/ABB.
\end{enumerate}

\subsection{Relationship to Existing Approaches}

Our approach is related to several strands of the density estimation literature.

\subsubsection*{Log-Spline Density Estimation}

The general framework of modelling $\log f(x) = S(x) - \log C$ where $S$ is
a spline has a long history (e.g., Stone 1990, Kooperberg and Stone 1991).
Our model is a specific instance with $S$ being a cubic spline on three
quantile knots, where the knot locations are determined by the quantile
regression parameters $a_Q$.

\subsubsection*{Gao and Hastie (2022, LinCDE)}

Gao and Hastie (2022) propose conditional density estimation via
$\log f(z|\mathbf{x}) = \sum_j \theta_j(\mathbf{x}) b_j(z) - c(\theta, \mathbf{x})$,
where $b_j$ are basis functions and $\theta_j(\mathbf{x})$ are
covariate-dependent coefficients.
Our model has this structure: $\mathbf{x} = \eta_{i,t-1}$,
$z = \eta_{it}$, and the basis functions are the cubic spline evaluated at
$z$ with knots at $Q_\tau(\eta_{i,t-1})$.
The coefficients $\theta_j(\eta_{i,t-1})$ are determined by the spline
values $s_1, s_2, s_3$ which are functions of the quantile knots
(themselves functions of $\eta_{i,t-1}$ through the Hermite basis).

\subsubsection*{Direct Extension of AB/ABB}

Unlike the general basis-expansion approach, our model is a \emph{direct}
smoothing of the AB/ABB quantile framework:
\begin{itemize}
\item The quantile knots $Q_{\tau_\ell}(\eta_{i,t-1})$ remain the primary
  parameters, just as in ABB.
\item The density between knots is determined by a cubic spline
  (instead of being piecewise uniform), with the constraint that each
  inter-quantile segment integrates to $\tau_{\ell+1} - \tau_\ell$.
\item The tails are quadratic in $\log f$ (Gaussian-like decay) instead of
  linear (exponential), with curvature parameter $\kappa_{\text{mean}}$
  determined from the inter-quartile range: $\kappa = -1.349^2 / \text{IQR}^2$.
\item The $C^2$ continuity at the knots is enforced analytically, not by
  ad hoc blending.
\end{itemize}

\subsection{Model Specification}

\subsubsection*{Log-Density}

On each segment, the log-density is a cubic polynomial.
For three quantile knots $t_1 < t_2 < t_3$ with log-density values
$s_1, s_2, s_3$ (where $s_2 = 0$ by normalisation), the cubic spline
$S(x)$ satisfies:
\begin{itemize}
\item $S(t_\ell) = s_\ell$ for $\ell = 1, 2, 3$ (interpolation),
\item $S''(t_1^+) = S''(t_1^-)$, $S''(t_3^-) = S''(t_3^+)$ ($C^2$ continuity),
\item the interior second derivative $M_2 = S''(t_2)$ is determined by
  the $C^1$ condition at $t_2$.
\end{itemize}

\subsubsection*{Tails}

The tails are quadratic in $\log f$:
\begin{align}
\log f(x) &= s_1 + \beta_L(x - t_1) + \tfrac{1}{2}\kappa_1(x - t_1)^2,
  \quad x < t_1, \\
\log f(x) &= s_3 + \beta_R(x - t_3) + \tfrac{1}{2}\kappa_3(x - t_3)^2,
  \quad x > t_3,
\end{align}
where $\beta_L = S'(t_1^+)$, $\beta_R = S'(t_3^-)$ (from $C^1$),
and $\kappa_1, \kappa_3 < 0$ ensure integrability.
The curvature parameters are linked via $\kappa_{\text{mean}} = (\kappa_1 + \kappa_3)/2$,
with the solver finding $\delta = (\kappa_3 - \kappa_1)/2$ to satisfy
the mass constraints.

\subsubsection*{Mass Constraints and the Newton Solver}

The three mass constraints
$\int_{\text{segment}_\ell} f(x)\, dx = 0.25$ for $\ell = 1, 2, 3$
determine $(s_1, s_3, \delta)$ given $(t_1, t_2, t_3, \kappa_{\text{mean}})$.
These are solved by a $3 \times 3$ Newton iteration.
The interior integrals $\int \exp(\text{cubic})$ are computed exactly via a
convergent Taylor series (typically 44 terms).

\subsubsection*{Profiled Parameterisation}

To match the parameter count of QR and eliminate weakly identified tail
parameters, we profile out $\kappa_{\text{mean}}$ using the IQR:
$\kappa_{\text{mean}} = -1.349^2 / (t_3 - t_1)^2$.
The resulting model has 14 parameters:
9 (transition quantile knots) + 3 (initial quantile knots) + 2 (epsilon quantile knots).

%==========================================================================
\section{Performance of MLE vs.\ QR in the Smooth ABB Model}
\label{sec:mc_smooth}
%==========================================================================

\subsection{Simulation Design}

We generate data from the smooth ABB model with $K = 2$, $\rho = 0.8$,
$\sigma_v = 0.5$, $\sigma_\varepsilon = 0.3$, $\sigma_{\eta_1} = 1.0$,
and a small quadratic heterogeneity $d_2 = 0.005$ in the transition quantile
gaps.
Data are generated using rejection sampling from the exact model density
(not grid-based CDF inversion, which we found introduces systematic bias;
see Section~\ref{sec:identification}).
Sample size $N = 300$, with 20 Monte Carlo replications.

MLE uses the profiled parameterisation (14 parameters, $\kappa$ from IQR)
with numerical gradient (adaptive step $h_j = \max(10^{-5}, 10^{-4}|v_j|)$)
and L-BFGS with BackTracking line search.
QR uses stochastic EM with 30 iterations and 10 FFBS draws per iteration,
estimating $\kappa$ from the IQR of the estimated quantiles.

\subsection{Results: $T = 3$}

\begin{table}[H]
\centering
\caption{Profiled MLE vs.\ QR, $T = 3$, $N = 300$, 20 replications.}
\label{tab:profiled_T3}
\begin{tabular}{lcccccccc}
\toprule
 & True & ML bias & ML std & ML RMSE & QR bias & QR std & QR RMSE & QR/ML \\
\midrule
$\rho$       & 0.800 & $-$0.001 & 0.036 & 0.036 & $-$0.006 & 0.030 & 0.030 & 0.86 \\
$a_{\varepsilon,3}$ & 0.202 & $-$0.008 & 0.045 & 0.046 & $+$0.002 & 0.023 & 0.023 & 0.50 \\
$a_{\text{init}}[1]$ & $-$0.674 & $-$0.026 & 0.101 & 0.105 & $-$0.010 & 0.081 & 0.081 & 0.78 \\
$a_{\text{init}}[2]$ & 0.000 & $-$0.032 & 0.090 & 0.096 & $-$0.035 & 0.082 & 0.089 & 0.93 \\
$a_{\text{init}}[3]$ & 0.674 & $-$0.038 & 0.074 & 0.084 & $-$0.042 & 0.088 & 0.098 & \textbf{1.17} \\
\bottomrule
\end{tabular}
\end{table}

\noindent
At $T = 3$, QR dominates MLE for $\rho$ (QR/ML $= 0.86$) and
$a_{\varepsilon,3}$ (QR/ML $= 0.50$), while MLE is better for $a_{\text{init}}[3]$
(QR/ML $= 1.17$).

\subsection{Results: $T = 4$}

\begin{table}[H]
\centering
\caption{Profiled MLE vs.\ QR, $T = 4$, $N = 300$, 20 replications.}
\label{tab:profiled_T4}
\begin{tabular}{lcccccccc}
\toprule
 & True & ML bias & ML std & ML RMSE & QR bias & QR std & QR RMSE & QR/ML \\
\midrule
$\rho$       & 0.800 & $+$0.002 & 0.029 & 0.029 & $-$0.000 & 0.033 & 0.033 & \textbf{1.16} \\
$a_{\varepsilon,3}$ & 0.202 & $-$0.012 & 0.041 & 0.043 & $-$0.001 & 0.015 & 0.015 & 0.35 \\
$a_{\text{init}}[1]$ & $-$0.674 & $-$0.012 & 0.077 & 0.078 & $+$0.002 & 0.070 & 0.070 & 0.89 \\
$a_{\text{init}}[2]$ & 0.000 & $-$0.013 & 0.078 & 0.079 & $-$0.022 & 0.071 & 0.074 & 0.94 \\
$a_{\text{init}}[3]$ & 0.674 & $-$0.026 & 0.070 & 0.075 & $-$0.048 & 0.076 & 0.090 & \textbf{1.20} \\
\bottomrule
\end{tabular}
\end{table}

\noindent
At $T = 4$, MLE now dominates QR for $\rho$ (QR/ML $= 1.16$) and continues
to dominate for $a_{\text{init}}[3]$ (QR/ML $= 1.20$).
However, QR remains substantially better for $a_{\varepsilon,3}$
(QR/ML $= 0.35$), and the gap \emph{widens} with more periods.

\textcolor{myred}{[NOTE: These results are preliminary (20 replications, $N = 300$).
They will be replaced by comprehensive simulation results from HPC
with larger $N$, more replications, and additional configurations including
$G \in \{201, 501, 1001\}$.]}

%==========================================================================
\section{Identification and Information Analysis}
\label{sec:identification}
%==========================================================================

\subsection{Local Identification: Information Matrix}

The observed information matrix $I_Y$ characterises the local precision of
the MLE.
By the Louis (1982) decomposition,
\begin{equation}\label{eq:louis}
I_Y = E[B(X,\theta) \mid Y] - \text{Var}[S(X,\theta) \mid Y]
    = I_X - I_{X \mid Y},
\end{equation}
where $X = (\eta_1, \ldots, \eta_T)$ is the complete data,
$I_X = E[B \mid Y]$ is the complete-data information (block diagonal in
$\theta_Q$, $\theta_{\text{init}}$, $\theta_\varepsilon$), and
$I_{X \mid Y} = \text{Var}[S \mid Y]$ is the missing information from
not observing $\eta$.

\subsubsection*{Block Diagonal Structure}

The complete-data log-likelihood decomposes as
\[
\lambda(X, \theta) = \log f_{\eta_1}(\eta_1; \theta_{\text{init}})
+ \sum_{t=2}^T \log f(\eta_t \mid \eta_{t-1}; \theta_Q)
+ \sum_{t=1}^T \log f_\varepsilon(y_t - \eta_t; \theta_\varepsilon).
\]
Since the three components depend on disjoint parameter blocks, the
complete-data information $I_X = E[B \mid Y]$ is \emph{block diagonal}.
The cross-parameter correlations in $I_Y$ arise entirely from the missing
information $I_{X \mid Y} = \text{Var}[S \mid Y]$, which has off-diagonal
blocks because conditioning on $Y$ induces correlations between the
scores of different components through the shared latent $\eta$.

\subsubsection*{Eigenvalue Analysis}

We compute $I_Y$ via the SML Hessian ($R = 2000$, $N = 300$, $T = 3$)
evaluated at the true parameter for the full (17-parameter) model.
The condition number is $2{,}810$, with all eigenvalues positive.
Table~\ref{tab:eigenvalues} reports the eigenvalues and their principal
directions.

\begin{table}[H]
\centering
\caption{Eigenvalues of $I_Y$ (SML Hessian, $R = 2000$, $N = 300$, $T = 3$),
sorted from weakest to strongest identified direction.}
\label{tab:eigenvalues}
\begin{tabular}{rll}
\toprule
Rank & Eigenvalue & Principal direction \\
\midrule
1 & $5.0 \times 10^7$ & $\delta_{L1}$ (left gap intercept) \\
2 & $5.0 \times 10^7$ & $\delta_{R1}$ (right gap intercept) \\
3 & $2.4 \times 10^8$ & $\log(-M_\varepsilon)$ (eps curvature) \\
4--5 & $5.5 \times 10^8$ & $\log(-a_{\varepsilon,1})$, $\log(a_{\varepsilon,3})$ \\
6 & $7.3 \times 10^8$ & $\log(-M_Q)$ (transition curvature) \\
7--8 & $1.2 \times 10^9$ & $\delta_{L2}$, $\delta_{R2}$ (gap slopes) \\
9 & $2.2 \times 10^9$ & $a_{Q,\text{med},1}$ ($\rho$ slope) \\
10--11 & $9.0 \times 10^9$ & Init gaps \\
12--17 & $1.4 \times 10^{11}$ & Intercepts, quadratic terms \\
\bottomrule
\end{tabular}
\end{table}

The weakest directions are the gap intercepts ($\delta_{L1}, \delta_{R1}$),
which control the IQR of the transition density at $\eta = 0$.
The curvature parameters ($M_\varepsilon, M_Q$) rank 3rd and 6th ---
substantially less well identified than the location parameters.
This is consistent with the simulation finding that MLE has high variance
for these parameters.

\subsection{Global Identification}

ABB establish nonparametric identification via the hidden Markov structure,
building on Hu and Schennach (2008) and Wilhelm (2015).
The key conditions are:
\begin{enumerate}
\item \textbf{Completeness}: the conditional expectation operators
  $\mathcal{L}_{y_2 \mid y_1}$ and $\mathcal{L}_{\eta_2 \mid y_1}$ are injective.
\item \textbf{Characteristic function non-vanishing}:
  $\mathcal{F}(s_1)(u) \neq 0$ for all $u \in \mathbb{R}$.
\end{enumerate}
These conditions ensure that $f_\varepsilon$ is identified via an ODE in
Fourier space (equation (S16) in the ABB supplement).

In our parametric model, all parameters are identified (all eigenvalues
of $I_Y$ are positive).
However, the identification of shape parameters (curvature, tail decay)
is achieved through \emph{deconvolution} --- recovering $f_\varepsilon$
from the convolution $f_y = f_\eta * f_\varepsilon$.
Deconvolution is an ill-posed inverse problem: high-frequency features
of $f_\varepsilon$ (such as tail curvature $M_\varepsilon$) are determined
by ratios of small Fourier coefficients, leading to slow convergence rates.

QR bypasses deconvolution entirely: it draws $\eta$ from the posterior
via FFBS and estimates $\varepsilon = y - \hat\eta$ directly.
This gives $\sqrt{N}$ rates for all parameters, whereas MLE's effective
rate for tail parameters is slower due to the ill-posed inversion.

\subsubsection*{Effect of More Time Periods}

Additional time periods provide more convolution observations, strengthening
identification.
Table~\ref{tab:T_comparison} summarises the QR/ML RMSE ratios.

\begin{table}[H]
\centering
\caption{QR/ML RMSE ratios by number of time periods (profiled MLE, $N = 300$).
Values $> 1$ indicate MLE is better.}
\label{tab:T_comparison}
\begin{tabular}{lccc}
\toprule
Parameter & $T = 3$ & $T = 4$ & Direction \\
\midrule
$\rho$ & 0.86 & \textbf{1.16} & MLE improves \\
$a_{\varepsilon,3}$ & 0.50 & 0.35 & QR advantage widens \\
$a_{\text{init}}[3]$ & \textbf{1.17} & \textbf{1.20} & MLE consistently better \\
\bottomrule
\end{tabular}
\end{table}

\noindent
MLE catches up for $\rho$ (the transition slope) as $T$ increases, but
the epsilon quantile $a_{\varepsilon,3}$ remains challenging.
This is consistent with the deconvolution interpretation: $\rho$ is
identified from the autocovariance structure (which improves with $T$),
while $a_{\varepsilon,3}$ requires separating $f_\varepsilon$ from the
convolution.

\subsection{Data Generation and Numerical Accuracy}

During the course of this investigation, we identified several sources of
numerical error that significantly affected MLE performance:

\begin{enumerate}
\item \textbf{Data generation}: The original \texttt{cspline\_draw} function
  used a 500-point trapezoidal CDF for inverse sampling, generating data
  from a density that did not exactly match the model.
  MLE exploited this misspecification, finding parameter values that fit
  the discretised data better than the true parameters.
  QR was unaffected because it estimates from residuals, not the density shape.
  Replacing with rejection sampling (exact draws from the model density)
  eliminated this bias.

\item \textbf{Gradient computation}: A fixed finite-difference step
  $h = 10^{-3}$ was inappropriate for parameters of different scales
  (e.g., $\log(-M_Q) \approx 1.4$ vs.\ gap $\delta \approx -5.3$).
  Replacing with adaptive steps $h_j = \max(10^{-5}, 10^{-4}|v_j|)$
  improved gradient accuracy.

\item \textbf{Tail parameterisation}: When the curvature parameters
  $M_Q, M_\varepsilon$ were freely estimated, MLE exhibited high variance
  (std $= 2.4$ for $M_Q$, truth $= -4.0$) and occasional NaN divergence
  (1/20 replications).
  Profiling $M$ from the IQR eliminated NaN and reduced variance by
  40\% for $M_Q$.
\end{enumerate}

%==========================================================================
\section{Conclusion}
\label{sec:conclusion}
%==========================================================================

This paper has studied the ABB nonlinear panel data framework from the
perspective of likelihood-based estimation.
We find that the original piecewise-uniform density creates fundamental
obstacles for MLE, and propose a smooth $C^2$ cubic spline extension that
preserves the quantile interpretation while enabling gradient-based optimisation.

Under correct specification of the smooth model, MLE and QR have complementary
strengths:
MLE is more efficient for the transition slope $\rho$ (especially with $T \geq 4$)
and the upper tail of the initial distribution, while QR is more efficient for
the transitory shock parameters, particularly the epsilon quantile
$a_{\varepsilon,3}$.
The QR advantage for epsilon parameters is consistent with the deconvolution
interpretation: QR bypasses the ill-posed inverse problem of separating
$f_\varepsilon$ from the convolution $f_y = f_\eta * f_\varepsilon$ by
working directly with $\eta$ draws from the posterior.

Several practical findings emerge:
(i) data generation must use exact sampling (rejection sampling) rather than
grid-based CDF inversion;
(ii) the tail curvature parameters $M_Q, M_\varepsilon$ should be profiled
out (estimated from the IQR) rather than freely optimised;
(iii) analytical gradients, while requiring substantial implementation effort,
eliminate NaN divergence and reduce bias.

\textcolor{myred}{[TODO: Add discussion of misspecified density comparison
(Part 2 of the research plan) and the Louis decomposition results once the
computation is corrected.]}

\end{document}
